\section{The Translation of Sincerity}
\label{sec:translation}

\noindent
Every chapter to this point has worked on the assumption that once genuine trust is produced, it is received\,---\,that a real self-commitment, supplied to another, lands as the ground it is. This section removes that assumption, and in removing it opens the problem that is, for the predicament of the Prelude, the most pressing of all. \emc{Sincerity has no universal language.} There is no form of sincerity legible to everyone; sincerity must be encoded to be received, and the language of the encoding depends upon the epistemology of the one who is to receive it. The gravity of the betrothal document, the moon, the poem, the philosophical seriousness\,---\,these encode sincerity in a form legible to one who reads character directly, who takes the bearing and the solemnity and the cost of the gesture as themselves the signal. To a receiver of another epistemology\,---\,one who holds that what is the case is the case and asks for what can be observed and verified\,---\,that same sincerity, in that same encoding, is not evidence at all. The same real sincerity, passed through two different decoders, yields utterly different degrees of credibility; and the difference is not in the sincerity, which is one and real, but in the mismatch between the sender's language of encoding and the receiver's language of decoding.

This is the point at which the paper's question undergoes its decisive turn. The question has been, through the mechanism and its consequences, an \emph{ontological} one: what is genuine trust, and how is it produced. It now becomes a \emph{translational} one: how is genuine trust, once produced, transmitted across the border between incommensurable epistemologies. And this translational question is the sharpest form, in the field of trust, of the master theme the series has carried throughout\,---\,that there is no metalanguage, that the Other does not exist as a guarantor standing above all positions. There is no neutral universal language of sincerity, no God's-eye dictionary in which the sincerity encoded in one epistemology can be looked up and re-expressed in another with its value preserved. The problem of trust, fully stated, is therefore not only how to be genuine but how to be \emph{read} as genuine by one whose language of reading is not one's own; and this section is its theory.

\subsection{The epistemology of translation}
\label{sub:trans-epistemology}

Before one can translate sincerity into a language a receiver can read, one must know what languages there are\,---\,what grammars of trust a receiver might be using, by which a profession of sincerity is parsed into a ground for expectation or rejected as no ground at all. This subsection sets out a genealogy of these grammars. It is offered as adequate to the practical purpose rather than as exhaustive, and its precision matters, for the effectiveness of any translation depends on the accuracy with which the target grammar is identified.

\subsubsection{A genealogy of trust-grammars}
\label{sub:trans-grammars}

There are at least four grammars of trust, irreducible to one another, each parsing a different kind of thing as a ground for expectation, and each with its characteristic reach and its characteristic blind spot.

The first is the \emph{evidentialist} or positivist grammar. It anchors trust in observable evidence and accomplished fact\,---\,the track record, the thing already done, the stable and verifiable pattern. To one who uses this grammar, a ground for trust is something that has shown itself, and a profession that has not yet shown itself in fact is not yet a ground. Its reach is into the demonstrated; its blind spot is the sincere newcomer who has, as yet, no record to show, and whom it must, by its own logic, misjudge as untrustworthy not because he is but because he has not yet had occasion to demonstrate otherwise. This is the grammar of Holmes's bad-man observer (\xref{sub:jur-holmes}) turned to the reading of others, and it is, the paper has reason to think, the grammar most relied upon by the family of the Prelude.

The second is the \emph{deontological} or characterological grammar. It anchors trust in a direct, intuitive reading of the other's character and sincerity\,---\,in the bearing, the manner, the quality of the commitment, taken as immediate signals of an inner state. To one who uses this grammar, the gravity of a vow, the seriousness of a gesture, the felt sincerity of a person, are themselves grounds, read off without the detour through accomplished fact. Its reach is into the inner and the present; its danger is that it reads too much, that it can be deceived by one who has learned to perform the signs of a sincerity he does not have. This is, the paper has reason to think, the grammar most native to the one who would encode his sincerity in red silk and the moon.

The third is the \emph{inductive} or relational-historical grammar. It anchors trust in the accumulated pattern of a shared history with this particular person\,---\,not in evidence as such, nor in a reading of character as such, but in the texture of what has passed between us, the record of how this one has been with me over time. Its reach is into the known particular; its blind spot is the wholly new relation, for which no shared history yet exists.

The fourth is the \emph{institutional} or endorsement grammar. It anchors trust in the vouching of a trusted third party\,---\,the matchmaker, the elder, the credential, the recognition of a community, the recommendation along the chain of \xref{sec:relational}. Its reach is into whatever the trusted institution can underwrite; its blind spot is the condition, diagnosed throughout this paper, in which the institutions of endorsement have themselves come under the crisis of the symbol and can no longer reliably underwrite anything.

These four are not a ranking, with one the true grammar and the rest its approximations. They are incommensurable epistemic positions, each parsing as a ground something the others do not, none reducible to the rest\,---\,and their irreducible plurality is one more vindication, in the field of trust, of the polyphonic law the series has insisted upon. The question whether the genealogy is fine enough is left open, and one refinement in particular is marked as undecided: whether the evidentialist grammar should be subdivided, between one that asks for accomplished fact and one that asks for demonstrable capacity, since these have different practical implications for what a translation into that grammar would have to supply.

\subsubsection{The notion of a grammar-configuration}
\label{sub:trans-configuration}

No person uses a single grammar purely. Each uses a \emph{configuration}\,---\,a weighting, mostly unreflective, of the several grammars, in which one may dominate without excluding the others. The family of the Prelude is best understood not as purely evidentialist but as running a configuration in which the evidentialist grammar dominates, the institutional has some weight, and the deontological is heavily discounted\,---\,so that a profession of character lands lightly and an accomplished fact lands hard. The notion of the configuration is what makes the genealogy a practical instrument rather than a typology: one does not translate one's sincerity into \emph{the} evidentialist grammar in the abstract, but into the particular configuration a particular receiver runs, and the diagnosis of that configuration is the first task of any actual translation. The unit of translation is the configuration, and the genealogy is the alphabet from which configurations are spelled.

That there can be no single correct translation of a sincerity from one grammar into another\,---\,that translation is, in principle, indeterminate\,---\,is the deep epistemological fact underlying all of this, and it is Quine's. The indeterminacy of translation~\citep{quine1960} is not a defect to be overcome but the very form the absence of a metalanguage takes in the theory of translation: there is no fact of the matter, fixed independently of every scheme, as to the uniquely correct rendering of a meaning from one language into another, and so no neutral standard against which a translation of sincerity could be certified as \emph{the} correct one. Translation is underdetermined all the way down, and the practical art the next subsection describes is an art practised under that permanent indeterminacy, not a technique for escaping it.

\subsection{The praxis of translation}
\label{sub:trans-praxis}

How, then, is a sincerity encoded in one grammar to be rendered legible in another, without distorting it, and without sliding into the manipulation that \xref{sec:justice} marked off? There are three paths, of ascending difficulty and depth, and each runs close to the ethical line, for each is, in its outward form, hard to distinguish from a counterfeit. The subsection sets them out, and then states the line that separates all three, done genuinely, from their forged twins.

\subsubsection{Projection}
\label{sub:trans-projection}

The first path is \emph{projection}: letting a sincerity that is legible in one's own grammar cast a shadow in the receiver's. One does not disguise the deontologically-legible sincerity as evidence, which would be forgery; one allows the real sincerity to leave, in the symbolic and observable register the receiver reads, the traces it naturally casts\,---\,an accumulation of consistent action, a constancy held across time, conduct that the evidentialist can read as evidence because it really is the observable shadow of a real inner state. The evidentialist infers back, from the shadow, to the substance. This is the safest path and the slowest, for it depends on time: the shadow accumulates only as the conduct accumulates, and cannot be hurried. It is the practical translation of the principle the series has carried as the maxim that what is held to over long enough comes to show itself\,---\,not that time creates the coupling-trust, which is real from the first, but that time lets the real cast, in the register the evidentialist reads, enough of a shadow to be recognised. Projection is the rendering of a real inner sincerity into the language of accomplished fact, by letting it accomplish the facts that are its true shadow.

\subsubsection{Multiple encoding}
\label{sub:trans-multiple}

The second path is \emph{multiple encoding}: encoding the one sincerity, at once, in two grammars, on two tracks running in parallel. To the one who reads character\,---\,and to oneself\,---\,the vow, the gravity, the poem; to the one who reads evidence, the stable and visible facts, the observable constancy, the demonstrated reliability. The two encodings are of one sincerity, offered each in the grammar that can receive it. The difficulty of this path is not technical but ethical, and it is severe: multiple encoding degrades into mere placation unless one genuinely concedes the legitimacy of the other's grammar. If one encodes in the evidentialist grammar while privately holding that to ask for evidence is shallow, that the deontological reading is the real one and the evidentialist a concession to the obtuse, then the second encoding is not a translation but a condescension, and a receiver of any acuity will read the condescension beneath it. Multiple encoding done genuinely requires that one hold both grammars as legitimate ways of arriving at trust\,---\,a requirement that connects it directly to the justice of presentation taken up in \xref{sec:praxis}, where the active rendering of one's sincerity in the other's grammar is shown to be, when the substance is real, not a strategy but a form of respect.

\subsubsection{Mutual translation}
\label{sub:trans-mutual}

The third path is the deepest and the hardest, and it is the one the section most wants to establish, for it is the only one that issues in a trust that is neither projected nor doubly-broadcast but jointly made. In \emph{mutual translation}, both parties move a step toward the other's language: one learns to speak in evidence, and the other learns to read the signals of character; and what is produced is not the persuasion of one by the other but the growing, in the interaction itself, of a small shared grammar that did not exist before. This is real coupling at the epistemological level\,---\,trust that is not one party's conviction transmitted to the other but a thing generated between them, in the continued interaction, out of two grammars that have each moved toward the other and grown, between them, a patch of common language. It does not pre-exist the relation and wait to be transmitted; it is generated, in the mutual movement, as the relation's own.

That this is not a philosopher's fancy but a real and learnable craft is shown by the existence of a mature profession that does exactly this. The experienced consultant, in any field, who is brought before a client he does not yet understand, faces the structural twin of the no-metalanguage problem: the client comes with a language, a situation, a set of terms and unspoken assumptions and a real difficulty, that the consultant cannot at first read; and the consultant can neither impose his own framework upon it, which is bad consulting, nor pretend to understand what he does not. The real work is to create, in the interaction, with this particular client, a common language\,---\,to learn the client's terms and let the client learn his, until a shared grammar exists between them that existed before in neither. This is mutual translation, and that consulting is a mature profession, with methods and training and tests of competence, is the proof that mutual translation can be done, systematically and well, and is not the unoperable ideal it might otherwise seem\,---\,the very objection this hardest path most invites. The empirical literature on the outcomes of such work lends this support: across the helping and advisory professions, the strongest predictor of a good outcome is found to be not the technique or the school but the quality of the working relationship, the working alliance\,---\,which is precisely the establishment of the shared understanding that mutual translation produces. The point answers, in part, the call the foundational paper made for empirical follow-up~\citep{wanhong_grb}: there is a body of evidence that the generation of a shared grammar is real, learnable, and the operative factor in success.

Two features of the consultant's craft sharpen the account. The first is an asymmetry. Mutual translation need not be symmetrical: the consultant, as the active party who bears the professional responsibility, moves further toward the client's language than the client moves toward his. The party who actively seeks to establish trust\,---\,the one in the position of the Prelude's suitor\,---\,may rightly have to move further toward the other's grammar, and this is not a mark of his grammar's inferiority but of his bearing the greater share of the labour of translation, as the active and responsible party. This asymmetry is the concrete form of what \xref{sec:praxis} will call epistemic hospitality: the moving-toward-the-other, the taking-on of the greater share of the work of being understood. The second feature names the philosophical resource that underwrites the path. What mutual translation accomplishes is the fusion of two horizons into a shared one that neither held alone\,---\,Gadamer's \emph{Horizontverschmelzung}~\citep{gadamer2004}, the fusion of horizons, which is the standing philosophical name for the generation of a common understanding between parties who began with none, and which gives mutual translation its grounding in the hermeneutic tradition.

\subsubsection{The line that separates translation from manipulation}
\label{sub:trans-line}

All three paths run close to the ethical line, and the line must be drawn as firmly as it can be, for it is the line on which the whole praxis turns and it is not, on the surface, where one might expect. To render one's sincerity legible in another's grammar (translation) and to manufacture a false signal legible in another's grammar (manipulation) may, as outward acts, look identical\,---\,the projected shadow and the forged evidence, the second encoding and the placating performance, can present the same surface. The difference lies wholly in whether what is being translated \emph{really exists}. To project the real shadow of a real sincerity is translation; to project the false shadow of a sincerity that is not there is manipulation (\xref{sec:justice}). The criterion of genuine translation is that the thing it translates must really exist and be capable of real fulfilment: it carries a real substance across the border between grammars, where manipulation manufactures, in the target grammar, the appearance of a substance that is absent. Translation takes the other's grammar as a legitimate other to be honestly addressed; manipulation takes the other's grammar as a lock to be picked. The line is not between act and act, which can match, but between a translated substance that is real and a manufactured signal that is empty\,---\,and it is, once more, the line of \xref{sec:justice}, the line between honouring the other and rendering him fuel.

\subsection{Dialogue: the convergence of three identities}
\label{sub:trans-dialogue}

The deepest of the three paths, mutual translation, turns out to be a thing the philosophical tradition has known under another name, and the recognition gives the section its close. Mutual translation \emph{is dialogue}, and dialogue is, at once, three things that prove to be one.

It is, first, the practical form of mutual translation: the consultant's craft, the back-and-forth in which a shared language is grown between parties who began without one. It is, second, the ancient method by which philosophy has always held that certain truths are generated\,---\,the Socratic \emph{dialektik\=e}, in which truth is not stated by one party and received by the other but generated between them, in the going and returning of question and answer, because some truths can be reached in no other way than by two minds moving toward each other. And it is, third, the formal law of the series itself\,---\,the polyphony of the plural conclusions, the refusal of a governing voice, the insistence that the truth of the matter lives not in any single framework but in the dialogue among them. These three\,---\,the consultant's craft, the Socratic method, the series' polyphonic form\,---\,are one thing, and their being one is why mutual translation is the deepest of the three paths. It is the deepest because it \emph{is} dialogue; and dialogue is, in a world without a metalanguage, the only way that truth and trust can be generated between heterogeneous subjects at all. Where there is no God's-eye dictionary to translate one grammar into another with the value preserved, the value is not transmitted but \emph{generated}, between the parties, in the dialogue in which each moves toward the other and a shared grammar grows. \emc{Mutual translation is dialogue: the generating, between two who share no common language, of a common language, together.}

And this returns the section to the poetics the series has carried throughout, for the translation of sincerity, so understood, is itself a poetic act in the precise sense the series has given the term. To translate one's sincerity into a grammar another can read is not to settle a meaning and hand it over closed; it is to invite the other to walk a path of interpretation, to offer a form the other must traverse rather than merely receive, to write, for the other, a poem he can read. The poetic refuses settlement and opens a path to be walked; and the translation of sincerity, done genuinely, refuses to settle the other's trust by manufacture and instead opens, between the two grammars, a path along which a shared understanding may be walked into being. The translation of sincerity is the writing, for the one whose language is not yours, of a poem in a language he can read\,---\,and the paper turns now from the technique of this writing to its practice, at the scale of institutions and then at the scale of a life.
